%! AP Statistics — Unit 9, Lesson 9.3 — Exit Ticket (Answer Key)
\documentclass[10pt]{article}
\usepackage{apstats-article}
\usepackage{apstats-key}

\begin{document}

\pageheader{Unit 9, Lesson 9.3}{Exit Ticket --- Answer Key}
\namedateperiod

\vspace{0.1in}
A researcher fits a regression line relating a person's age (years, $x$) to their reading
speed (words per minute, $y$) for adults enrolled in a library program. From a random sample
of 36 adults, computer output gives a 95\% confidence interval for the slope:
$(-1.4,\ -0.2)$ wpm per year.

\begin{enumerate}

  \item Write a complete interpretation of this confidence interval in context.

        \ans{We are 95\% confident that the true slope of the population regression line
        relating reading speed to age for adults in this library program is between $-1.4$
        and $-0.2$ words per minute per year.}

  \item Does this CI provide evidence that $\beta \ne 0$? Justify your answer.

        \ans{Yes. Since 0 is not contained in the interval $(-1.4,\ -0.2)$, this provides
        evidence that $\beta \ne 0$; there is a negative linear relationship between age and
        reading speed for adults in this library program.}

  \item A reading specialist claims that age decreases reading speed by exactly
        1 word per minute per year ($\beta = -1$). Is this claim consistent with the
        confidence interval? Explain.

        \ans{Yes. Since $-1$ is contained in the interval $(-1.4,\ -0.2)$, the specialist's
        claim that $\beta = -1$ is consistent with the data.}

\end{enumerate}

\begin{teachernote}
Students should write ``consistent with'' (not ``proves'' or ``confirms'') when the claimed
value falls inside the interval. Full credit on Question 1 requires all three elements:
confidence level, both endpoints with units, and a reference to the population (not just the
sample). On Question 2, students must identify that 0 is outside the interval to earn credit.
\end{teachernote}

\end{document}
